\documentclass[a4paper,10pt]{report}
\usepackage{geometry}\geometry{a4paper,top=3.5cm,bottom=3.5cm,%
left=2.5cm,right=2.5cm,heightrounded,bindingoffset=0mm}
\usepackage[T1]{fontenc}
\usepackage[utf8]{inputenc}
\usepackage[italian]{babel}
\usepackage{graphicx}
\usepackage[export]{adjustbox}%Per il Frame attrono le immagini e il valign
\usepackage{subfig}
\usepackage{amsmath,amsfonts,amssymb,braket,mathrsfs}
\usepackage{float}
\usepackage{tabularx,booktabs}
\usepackage{hyperref}
\usepackage{epsfig}
\usepackage{pdfpages} %Per gli allegati
%\usepackage{minipage}
\usepackage[output-decimal-marker={,}]{siunitx}
\usepackage{tikz}
\usepackage{pgfplots,pgfplotstable}
%\pgfplotsset{compat=1.15} %indica la versione da utilizzare per pgfplot
\usetikzlibrary{patterns} % per il tratteggio
\usepgfplotslibrary{groupplots}
\pgfplotsset{compat=newest}
\usepackage{stanli}
\usepackage{xspace}% per lo spazio intelligente
\newcommand{\e}{\`E\xspace}  %E'
\usepackage{titlesec} % per formato custom dei titoli dei capitoli
%\usepackage{sideways}%%%
% redefinizione del formato del titolo del capitolo
      % da formato
      %   Capitolo X
      %   Titolo capitolo
      %   a formato
      %      Titolo capitolo 
	\titleformat{\chapter}
        {\normalfont\Huge\bfseries}{}{0em}{}
	\titlespacing*{\chapter}{0pt}{0in}{0.02in}
	\titlespacing*{\section}{0pt}{0.2in}{0.02in}
	\titlespacing*{\subsection}{0pt}{0.10in}{0.02in}
%serve per la didascalia di tabelle e figure:
\usepackage{caption}
\captionsetup{tableposition=top,figureposition=bottom,font=small}\captionsetup{format=hang,labelfont={bf,color=pantone186}} %didascalie a più righe allineate e il nome in grassetto
%non viene allineato a sinistra se la didascalia è corta una sola riga. PERCHé??
\usepackage{xcolor}
%serve per mettere il codice con lo sfondo grigio chiaro
\definecolor{pantone186}{RGB}{206, 17, 38} %il colore del logo UNITN
\definecolor{myGray}{gray}{0.5} %più basso più scuro è
\usepackage{listings} 
\lstset{basicstyle=\scriptsize\ttfamily,
backgroundcolor=\color{lightgray},%
boxpos=c,%
stringstyle=\itshape,		
lineskip=3pt,%
numbers=left,
numberstyle=\tiny,}
\usepackage{lscape}
\usepackage{multirow}
\usepackage{import}
%\usepackage{pythontex}
\begin{document}
%!TEX root = ../TesiTriennaleMeoliNicola.tex
\pagestyle{plain}
\thispagestyle{empty}
\begin{center}
  \begin{figure}[H]
    \centerline{\psfig{file=IMG/logo_unitn_black_centerNEW.eps,
                        width=0.8\textwidth,trim = 0 0.9cm 0 0.5cm}}
  \end{figure}
\textcolor{pantone186}{\noindent\rule{\textwidth}{.5pt}}

  \Large\textsc{Dipartimento di Ingegneria Civile, Ambientale e Meccanica\\}
  \Large{Corso di Laurea in Ingegneria Civile
  }

  \vspace{3.7 cm} 
  %\Large\textsc{Elaborato finale\\} 
  %\vspace{1 cm} 
  \Huge\textsc{Confronto energetico tra diverse soluzioni edilizie\\}
  
  \vspace{0.2 cm}
  \Large{\it{Analisi termo igrometrica di alcuni pacchetti strutturali composti da differenti materiali. }}


  \vspace{4 cm} 
  \begin{tabular*}{\textwidth}{ l @{\extracolsep{\fill}} r }
  \Large\textsc{Docenti} & \Large\textsc{Studente}\\
  \Large{Rossano Albatici}& \Large{Nicola Meoli 186100}\\
  
  	
  	
  \end{tabular*}

  \vspace{3.1cm} 
  \textcolor{pantone186}{\noindent\rule{\textwidth}{1pt}}
    
  \Large{Anno accademico 2020/21}
  
\end{center}


\tableofcontents
%\setcounter{page}{1}
%Tabelle e figure sulla stessa pagina:
%Le aggiunge all'indice. phantomsection serve per non far casini con hyperref
\clearpage
\begingroup
   %\let\cleardoublepage\relax  % book
    \let\clearpage\relax        % report
        \listoftables
        \phantomsection
        \addcontentsline{toc}{chapter}{Elenco delle tabelle}
        %
        \listoffigures
        \phantomsection
        \addcontentsline{toc}{chapter}{Elenco delle figure}
\endgroup
%

%!TEX root = ../TesiTriennaleMeoliNicola.tex
\chapter{Introduzione}
Sappiamo come nell'ambito della progettazione edilizia esistano un'infinità di soluzioni costruttive,ma quel è la migliore tra tutte? 
In un precedente lavoro di gruppo era stata analizzata la convenienza economica di diversi pacchetti costruttivi di parete, prendendo come parametri di confronto il costo iniziale, di manutenzione e il tempo di messa in opera. 
Per avere soluzioni simili era stato scelto come parametro fisso la resistenza termica ponendola pari a $R=\SI{4}{m^2K\per W}$.
Il confronto era tra pacchetti costituiti da materiali ben differenti tra loro; tra cui telaio in calcestruzzo armato, telaio in acciaio, muratura portante e parete in X-LAM portante.
A questi strati strutturali si erano poi aggiunti gli isolanti e le rifiniture interne ed esterne.

Per quanto riguarda lo strato strutturale si era arrivati alla conclusione che, dal punto di vista economico, conveniva di più la soluzione in calcelstruzzo armato. 
Perché sebbene richieda molto più tempo di posa in opera, essa è anche molto meno costosa. 
Al lato opposto si aveva la soluzione in X-LAM. 
Ovvero l'unica proposta scelta che prevedeva l'intera parete con una tipologia costruttiva prefabbricata.
Proprio per questo risultava la più rapida, e ciò permetteva di avere una durata del cantiere molto più breve rispetto le altre. 
In più richiedeva meno manutenzione.
Questo era un vantaggio economico nel lungo periodo, ma uno svantaggio enorme iniziale perché il costo di costruzione a causa solo del materiale, era molto più elevato.

(GRAFICO FINALE COSTI O DATI CON I COSTI)

Ora invece si vogliono prendere in esame soltanto due pacchetti e confrontarli non più dal punto di vista economico, ma da quello termo-igrometrico.
Si vuole cercare di capire, cioè, se nonostante i costi più elevati iniziali, ci siano dei vantaggi che facciano scegliere per il pacchetto composto da legno.

Nei capitoli successivi si entrerà più nel dettaglio, riportando un'analisi di Glaser effettuata e una ricerca accademica su questo materiale diametralmente opposto a tutti gli altri.

(DA AGGIUNGERE QUALCHE DATO DELLA VECCHIA RELAZIONE)
%!TEX root = ../TesiTriennaleMeoliNicola.tex
\chapter{Analisi statica}
\begin{align*}
    \Delta\vartheta &= \vartheta_{\text{int.}} - \vartheta_{\text{est.}} \\
    \vartheta^{(1)} &= \vartheta_{\text{int.}} \\
    \vartheta^{(k)} &= \vartheta_{\text{int.}} - \frac{\Delta\vartheta}{R_{tot}} \, \left(R_{\text{s,int.}} + \sum_{i=1}^{k}R^{(i)}\right) \\
    \vartheta^{(N)} &= \vartheta_{\text{est.}} \qquad \left[\SI{}{\celsius}\right]
\end{align*}

\begin{equation}
    p^{(k)}_{\text{SAT.}} =
    \begin{cases}
            610.5\exp\left(\frac{17.269\,\vartheta^{(k)}}{237.3 + \vartheta^{(k)}}\right) & \text{per $\vartheta^{(k)}\geq \SI{0}{\celsius}$}\\ 
            610.5\exp\left(\frac{21.875\,\vartheta^{(k)}}{265.5 + \vartheta^{(k)}}\right) & \text{per $\vartheta^{(k)} \leq  \SI{0}{\celsius}$}
    \end{cases}
    \qquad \left[\SI{}{\pascal}\right]
\end{equation}

\begin{align*}
    p^{(1)} &= p_{\text{int.}} = p^{\text{SAT.}}_{\text{int.}} \cdot UR_{\text{int.}}\\
    p^{(N)} &= p_{\text{est.}} = p^{\text{SAT.}}_{\text{est.}} \cdot UR_{\text{est.}}\\
    \Delta p &= p_{\text{int.}} - p_{\text{est.}} \\
    p^{(k)} &= p_{\text{int.}} - \frac{\Delta p}{Sd_{tot}} \, \left(\sum_{i=1}^{k} Sd^{(i)} \right) \qquad \left[\SI{}{\pascal}\right]
\end{align*}
\clearpage
\begin{landscape}
        \subsection*{Parete in muratura con isolante interno}
                \begin{minipage}[c]{0.4\linewidth}
                        \begin{table}[H]
\centering
\caption{Parete in muratura con isolante interno}
\begin{tabular}{lrrr}
\toprule
            Strati & Spessori & Conduttività  &    mu \\
\midrule
          Intonaco &    0,015 &                0,600 &   6,0 \\
  Isolante RockDuo &    0,140 &                0,035 &   1,0 \\
 Laterizio Poroton &    0,200 &                0,230 &  10,0 \\
          Intonaco &    0,015 &                0,900 &   8,0 \\
\bottomrule
\end{tabular}
\end{table}

                \end{minipage}
                \begin{minipage}[c]{0.6\linewidth}
                        % This file was created by tikzplotlib v0.9.4.
\begin{tikzpicture}

\definecolor{color0}{rgb}{1,0.647058823529412,0}

\begin{groupplot}[group style={group size=1 by 2,vertical sep=2.5cm}]
\nextgroupplot[
	 ticklabel style={ 
 		 /pgf/number format/fixed, 
 		 /pgf/number format/precision=5
		}, 
scaled ticks=false,
height=7cm,
legend cell align={left},
legend style={fill opacity=0.8, draw opacity=1, text opacity=1, draw=white!80!black},
minor xtick={},
minor ytick={},
tick align=outside,
tick pos=left,
width=\linewidth,
x grid style={white!69.0196078431373!black},
xlabel={Spessore parete (m)},
xmajorgrids,
xmin=-0.0185, xmax=0.3885,
xtick style={color=black},
xtick={0,0.015,0.155,0.355,0.37},
xticklabel style = {rotate=90.0},
y grid style={white!69.0196078431373!black},
ylabel={Temperature (°C)},
ymin=-6.01137676302392, ymax=20.5685707562642,
ytick style={color=black},
ytick={-10,-5,0,5,10,15,20,25}
]
\addplot [line width=1.5pt, green!50.1960784313725!black]
table {%
0 19.3603913235693
0.015 19.2373896550249
0.155 -0.442877312074838
0.355 -4.72119621796609
0.37 -4.803197330329
};
\addlegendentry{Temperatura}

\nextgroupplot[
	 ticklabel style={ 
 		 /pgf/number format/fixed, 
 		 /pgf/number format/precision=5
		}, 
scaled ticks=false,
height=7cm,
legend cell align={left},
legend style={fill opacity=0.8, draw opacity=1, text opacity=1, draw=white!80!black},
minor xtick={},
minor ytick={},
tick align=outside,
tick pos=left,
width=\linewidth,
x grid style={white!69.0196078431373!black},
xlabel={Spessore equivalente Sd (m)},
xmajorgrids,
xmin=-0.1175, xmax=2.4675,
xtick style={color=black},
xtick={0,0.09,0.23,2.23,2.35},
xticklabel style = {rotate=90.0},
y grid style={white!69.0196078431373!black},
ylabel={Pressione (Pa)},
ymin=266.815315862144, ymax=2340.26179765774,
ytick style={color=black},
ytick={250,500,750,1000,1250,1500,1750,2000,2250,2500}
]
\addplot [line width=1.5pt, color0]
table {%
0 1519.01824347152
0.09 1474.67101690856
0.23 1405.68644225507
2.23 420.192518633768
2.35 361.06288321649
};
\addlegendentry{Pressione}
\addplot [line width=1.5pt, red]
table {%
0 2246.0142303034
0.09 2228.8857975183
0.23 588.588856594798
2.23 410.858360238758
2.35 407.990247749789
};
\addlegendentry{Pressione Saturazione}
\end{groupplot}

\end{tikzpicture}
    
                \end{minipage}
        \subsection*{Parete in muratura con isolante esterno}
                \begin{minipage}[c]{0.4\linewidth}
                        \begin{table}[H]
\centering
\caption{Parete in muratura con isolante esterno}
\begin{tabular}{lrrr}
\toprule
            Strati & Spessori & Conduttività  &    mu \\
\midrule
          Intonaco &    0,015 &                0,600 &   6,0 \\
 Laterizio Poroton &    0,200 &                0,230 &  10,0 \\
  Isolante RockDuo &    0,140 &                0,035 &   1,0 \\
          Intonaco &    0,015 &                0,900 &   8,0 \\
\bottomrule
\end{tabular}
\end{table}

                \end{minipage}
                \begin{minipage}[c]{0.6\linewidth}
                        % This file was created by tikzplotlib v0.9.4.
\begin{tikzpicture}

\definecolor{color0}{rgb}{1,0.647058823529412,0}

\begin{groupplot}[group style={group size=1 by 2,vertical sep=2.5cm}]
\nextgroupplot[
	 ticklabel style={ 
 		 /pgf/number format/fixed, 
 		 /pgf/number format/precision=5
		}, 
scaled ticks=false,
height=7cm,
legend cell align={left},
legend style={fill opacity=0.8, draw opacity=1, text opacity=1, draw=white!80!black},
minor xtick={},
minor ytick={},
tick align=outside,
tick pos=left,
width=\linewidth,
x grid style={white!69.0196078431373!black},
xlabel={Spessore parete (m)},
xmajorgrids,
xmin=-0.0175, xmax=0.3675,
xtick style={color=black},
xtick={0,0.015,0.215,0.335,0.35},
xticklabel style = {rotate=90.0},
y grid style={white!69.0196078431373!black},
ylabel={Temperature (°C)},
ymin=-5.94820265073794, ymax=20.3881663325197,
ytick style={color=black},
ytick={-10,-5,0,5,10,15,20,25}
]
\addplot [line width=1.5pt, green!50.1960784313725!black]
table {%
0 19.1910586514625
0.015 19.035493007513
0.215 16.6873323441241
0.335 -4.64738454038109
0.35 -4.75109496968077
};
\addlegendentry{Temperatura}

\nextgroupplot[
	 ticklabel style={ 
 		 /pgf/number format/fixed, 
 		 /pgf/number format/precision=5
		}, 
scaled ticks=false,
height=7cm,
legend cell align={left},
legend style={fill opacity=0.8, draw opacity=1, text opacity=1, at={(0.03,0.03)}, anchor=south west, draw=white!80!black},
minor xtick={},
minor ytick={},
tick align=outside,
tick pos=left,
width=\linewidth,
x grid style={white!69.0196078431373!black},
xlabel={Spessore equivalente Sd (m)},
xmajorgrids,
xmin=-0.1165, xmax=2.4465,
xtick style={color=black},
xtick={0,0.09,2.09,2.21,2.33},
xticklabel style = {rotate=90.0},
y grid style={white!69.0196078431373!black},
ylabel={Pressione (Pa)},
ymin=267.992841074213, ymax=2315.53376820431,
ytick style={color=black},
ytick={250,500,750,1000,1250,1500,1750,2000,2250,2500}
]
\addplot [line width=1.5pt, color0]
table {%
0 1519.01824347152
0.09 1474.29035401961
2.09 480.337255088253
2.21 420.700069152371
2.33 361.06288321649
};
\addlegendentry{Pressione}
\addplot [line width=1.5pt, red]
table {%
0 2222.46372606204
0.09 2201.01860298296
2.09 1898.60822561279
2.21 413.455716135536
2.33 409.810487413109
};
\addlegendentry{Pressione Saturazione}
\end{groupplot}

\end{tikzpicture}
    
                \end{minipage}
        \subsection*{Perete in X-LAM con isolante bassa densità}
                \begin{minipage}[c]{0.4\linewidth}
                        \begin{table}[H]
\centering
\caption{Perete in X-LAM con isolante bassa densità}
\begin{tabular}{lrrr}
\toprule
                      Strati & Spessori & Conduttività  &    mu \\
\midrule
                  Gessofibra &    0,013 &                0,210 &   5,0 \\
                   X-LAM KLH &    0,096 &                0,130 &  25,0 \\
 Isolante bassa densità  &    0,120 &                0,038 &   1,0 \\
              Intonaco calce &    0,015 &                0,900 &  20,0 \\
\bottomrule
\end{tabular}
\end{table}

                \end{minipage}
                \begin{minipage}[c]{0.6\linewidth}
                        % This file was created by tikzplotlib v0.9.4.
\begin{tikzpicture}

\definecolor{color0}{rgb}{1,0.647058823529412,0}

\begin{groupplot}[group style={group size=1 by 2,vertical sep=2.5cm}]
\nextgroupplot[
	 ticklabel style={ 
 		 /pgf/number format/fixed, 
 		 /pgf/number format/precision=5
		}, 
scaled ticks=false,
height=7cm,
legend cell align={left},
legend style={fill opacity=0.8, draw opacity=1, text opacity=1, draw=white!80!black},
minor xtick={},
minor ytick={},
tick align=outside,
tick pos=left,
width=\linewidth,
x grid style={white!69.0196078431373!black},
xlabel={Spessore parete (m)},
xmajorgrids,
xmin=-0.0123, xmax=0.2583,
xtick style={color=black},
xtick={0,0.015,0.111,0.231,0.246},
xticklabel style = {rotate=90.0},
y grid style={white!69.0196078431373!black},
ylabel={Temperature (°C)},
ymin=-5.95814438534438, ymax=20.4165566468082,
ytick style={color=black},
ytick={-10,-5,0,5,10,15,20,25}
]
\addplot [line width=1.5pt, green!50.1960784313725!black]
table {%
0 19.2177065998922
0.015 18.7878750613714
0.111 14.3440782323564
0.231 -4.6590003127735
0.246 -4.75929433842835
};
\addlegendentry{Temperatura}

\nextgroupplot[
	 ticklabel style={ 
 		 /pgf/number format/fixed, 
 		 /pgf/number format/precision=5
		}, 
scaled ticks=false,
height=7cm,
legend cell align={left},
legend style={fill opacity=0.8, draw opacity=1, text opacity=1, draw=white!80!black},
minor xtick={},
minor ytick={},
tick align=outside,
tick pos=left,
width=\linewidth,
x grid style={white!69.0196078431373!black},
xlabel={Spessore equivalente Sd (m)},
xmajorgrids,
xmin=-0.14475, xmax=3.03975,
xtick style={color=black},
xtick={0,0.075,2.475,2.595,2.895},
xticklabel style = {rotate=90.0},
y grid style={white!69.0196078431373!black},
ylabel={Pressione (Pa)},
ymin=267.808252731272, ymax=2319.41012340606,
ytick style={color=black},
ytick={250,500,750,1000,1250,1500,1750,2000,2250,2500}
]
\addplot [line width=1.5pt, color0]
table {%
0 1519.01824347152
0.075 1489.0193999416
2.475 529.056406984059
2.595 481.058257336182
2.895 361.06288321649
};
\addlegendentry{Pressione}
\addplot [line width=1.5pt, red]
table {%
0 2226.15549292085
0.075 2167.25718410785
2.475 1633.75511312812
2.595 413.045980867543
2.895 409.52354638032
};
\addlegendentry{Pressione Saturazione}
\end{groupplot}

\end{tikzpicture}
    
                \end{minipage}
        \subsection*{Perete in X-LAM con isolante bassa densità lana di roccia}
                \begin{minipage}[c]{0.4\linewidth}
                        \begin{table}[H]
\centering
\caption{Perete in X-LAM con isolante bassa densità lana di roccia}
\begin{tabular}{lrrr}
\toprule
                Strati & Spessori & Conduttività &    mu \\
\midrule
            Gessofibra &    0,013 &                0,210 &   5,0 \\
             X-LAM KLH &    0,096 &                0,130 &  25,0 \\
 Isolante ventirockduo &    0,105 &                0,035 &   1,0 \\
        Intonaco calce &    0,015 &                0,900 &  20,0 \\
\bottomrule
\end{tabular}
\end{table}

                \end{minipage}
                \begin{minipage}[c]{0.6\linewidth}
                        % This file was created by tikzplotlib v0.9.4.
\begin{tikzpicture}

\definecolor{color0}{rgb}{1,0.647058823529412,0}

\begin{groupplot}[group style={group size=1 by 2,vertical sep=2.5cm}]
\nextgroupplot[
	 ticklabel style={ 
 		 /pgf/number format/fixed, 
 		 /pgf/number format/precision=5
		}, 
scaled ticks=false,
height=7cm,
legend cell align={left},
legend style={fill opacity=0.8, draw opacity=1, text opacity=1, draw=white!80!black},
minor xtick={},
minor ytick={},
tick align=outside,
tick pos=left,
width=\linewidth,
x grid style={white!69.0196078431373!black},
xlabel={Spessore parete (m)},
xmajorgrids,
xmin=-0.01145, xmax=0.24045,
xtick style={color=black},
xtick={0,0.013,0.109,0.214,0.229},
xticklabel style = {rotate=90.0},
y grid style={white!69.0196078431373!black},
ylabel={Temperature (°C)},
ymin=-5.94588914613306, ymax=20.3815597265862,
ytick style={color=black},
ytick={-10,-5,0,5,10,15,20,25}
]
\addplot [line width=1.5pt, green!50.1960784313725!black]
table {%
0 19.1848575050989
0.013 18.7966944122889
0.109 14.1662991749812
0.214 -4.64468147658159
0.229 -4.74918692464583
};
\addlegendentry{Temperatura}

\nextgroupplot[
	 ticklabel style={ 
 		 /pgf/number format/fixed, 
 		 /pgf/number format/precision=5
		}, 
scaled ticks=false,
height=7cm,
legend cell align={left},
legend style={fill opacity=0.8, draw opacity=1, text opacity=1, draw=white!80!black},
minor xtick={},
minor ytick={},
tick align=outside,
tick pos=left,
width=\linewidth,
x grid style={white!69.0196078431373!black},
xlabel={Spessore equivalente Sd (m)},
xmajorgrids,
xmin=-0.1435, xmax=3.0135,
xtick style={color=black},
xtick={0,0.065,2.465,2.57,2.87},
xticklabel style = {rotate=90.0},
y grid style={white!69.0196078431373!black},
ylabel={Pressione (Pa)},
ymin=268.035757521022, ymax=2314.63252282132,
ytick style={color=black},
ytick={250,500,750,1000,1250,1500,1750,2000,2250,2500}
]
\addplot [line width=1.5pt, color0]
table {%
0 1519.01824347152
0.065 1492.79277363996
2.465 524.46773370544
2.57 482.103513208305
2.87 361.06288321649
};
\addlegendentry{Pressione}
\addplot [line width=1.5pt, red]
table {%
0 2221.60539712585
0.065 2168.45182859209
2.465 1615.05395862344
2.57 413.551117171863
2.87 409.87728665048
};
\addlegendentry{Pressione Saturazione}
\end{groupplot}

\end{tikzpicture}
    
                \end{minipage}
        \subsection*{Perete in X-LAM con isolante alta densità fibra di legno}
                \begin{minipage}[c]{0.4\linewidth}
                        \begin{table}[H]
\centering
\caption{Perete in X-LAM con isolante alta densità fibra di legno}
\begin{tabular}{lrrr}
\toprule
                                    Strati & Spessori & Conduttività  &    mu \\
\midrule
                                Gessofibra &    0,013 &                0,210 &   5,0 \\
                                 X-LAM KLH &    0,096 &                0,130 &  25,0 \\
 Isolante alta densità &    0,130 &                0,043 &   5,0 \\
                            Intonaco calce &    0,015 &                0,900 &  20,0 \\
\bottomrule
\end{tabular}
\end{table}

                \end{minipage}
                \begin{minipage}[c]{0.6\linewidth}
                        % This file was created by tikzplotlib v0.9.4.
\begin{tikzpicture}

\definecolor{color0}{rgb}{1,0.647058823529412,0}

\begin{groupplot}[group style={group size=1 by 2,vertical sep=2.5cm}]
\nextgroupplot[
	 ticklabel style={ 
 		 /pgf/number format/fixed, 
 		 /pgf/number format/precision=5
		}, 
scaled ticks=false,
height=7cm,
legend cell align={left},
legend style={fill opacity=0.8, draw opacity=1, text opacity=1, draw=white!80!black},
minor xtick={},
minor ytick={},
tick align=outside,
tick pos=left,
width=\linewidth,
x grid style={white!69.0196078431373!black},
xlabel={Spessore parete (m)},
xmajorgrids,
xmin=-0.0127, xmax=0.2667,
xtick style={color=black},
xtick={0,0.013,0.109,0.239,0.254},
xticklabel style = {rotate=90.0},
y grid style={white!69.0196078431373!black},
ylabel={Temperature (°C)},
ymin=-5.94765269629741, ymax=20.3865958440658,
ytick style={color=black},
ytick={-10,-5,0,5,10,15,20,25}
]
\addplot [line width=1.5pt, green!50.1960784313725!black]
table {%
0 19.1895845467766
0.013 18.803672426194
0.109 14.2001290232679
0.239 -4.64674198192825
0.254 -4.75064139900817
};
\addlegendentry{Temperatura}

\nextgroupplot[
	 ticklabel style={ 
 		 /pgf/number format/fixed, 
 		 /pgf/number format/precision=5
		}, 
scaled ticks=false,
height=7cm,
legend cell align={left},
legend style={fill opacity=0.8, draw opacity=1, text opacity=1, draw=white!80!black},
minor xtick={},
minor ytick={},
tick align=outside,
tick pos=left,
width=\linewidth,
x grid style={white!69.0196078431373!black},
xlabel={Spessore equivalente Sd (m)},
xmajorgrids,
xmin=-0.17075, xmax=3.58575,
xtick style={color=black},
xtick={0,0.065,2.465,3.115,3.415},
xticklabel style = {rotate=90.0},
y grid style={white!69.0196078431373!black},
ylabel={Pressione (Pa)},
ymin=268.003044265693, ymax=2315.31950118321,
ytick style={color=black},
ytick={250,500,750,1000,1250,1500,1750,2000,2250,2500}
]
\addplot [line width=1.5pt, color0]
table {%
0 1519.01824347152
0.065 1496.97809752231
2.465 683.188093243512
3.115 462.786633751339
3.415 361.06288321649
};
\addlegendentry{Pressione}
\addplot [line width=1.5pt, red]
table {%
0 2222.25966223242
0.065 2169.39745901465
2.465 1618.59809228845
3.115 413.4783925652
3.415 409.826365687118
};
\addlegendentry{Pressione Saturazione}
\end{groupplot}

\end{tikzpicture}
    
                \end{minipage}
 \clearpage       
\end{landscape}
\clearpage
%!TEX root = ../TesiTriennaleMeoliNicola.tex
\chapter{Analisi dinamica}

Definizione di alcuni parametri utilizzati:
Profondità di penetrazione $\delta$ 
Massa superficiale
blabla



\begin{equation}
    \delta^{(k)} = \sqrt{\frac{\lambda^{(k)} \, T}{\pi \,\rho^{(k)} \,c^{(k)}}}
\end{equation}
La matrice di trasferimento, calcolata per ogni singolo strato $k$-esimo è definita come:
\begin{equation}
    \mathbf{Z}^{(k)} =
    \begin{bmatrix}  
        Z_{11}^{(k)} & Z_{12}^{(k)} \\
        Z_{21}^{(k)} & Z_{22}^{(k)}
    \end{bmatrix}
\end{equation}
dove le componenti della matrice $\mathbf{Z^{(k)}}$ sono

{\footnotesize
\begin{align*}
    Z_{11}^{(k)} &= Z_{22}^{(k)} = \cosh\left(\xi^{(k)}\right)\cos\left(\xi^{(k)}\right) + \mathrm{i} \left[\sinh\left(\xi^{(k)}\right)\sin\left(\xi^{(k)}\right)\right] \\
    Z_{12}^{(k)} &=-\frac{\delta^{(k)}}{2\lambda^{(k)}}\Bigl\{\sinh\left(\xi^{(k)}\right)\cos\left(\xi^{(k)}\right)+\cosh\left(\xi^{(k)}\right)\sin\left(\xi^{(k)}\right) + \mathrm{i}\left[\cosh\left(\xi^{(k)}\right)\sin\left(\xi^{(k)}\right)-\sinh\left(\xi^{(k)}\right)\cos\left(\xi^{(k)}\right)\right]  \Bigr\} \\
    Z_{21}^{(k)} &= \frac{\lambda^{(k)}}{\delta^{(k)}}\Bigl\{  \sinh\left(\xi^{(k)}\right)\cos\left(\xi^{(k)}\right)-\cosh\left(\xi^{(k)}\right)\sin\left(\xi^{(k)}\right) + \mathrm{i} \left[\sinh\left(\xi^{(k)}\right)\cos\left(\xi^{(k)}\right)+\cosh\left(\xi^{(k)}\right)\sin\left(\xi^{(k)}\right)\right] \Bigr\}
\end{align*}
}%
e con $\xi^{(k)}$ si intende il rapporto $\dfrac{d^{(k)}}{\delta^{(k)}}$ calcolato per ogni strato.

Moltiplicando le matrici di trasferimento di ogni singolo strato si ottiene la matrice di trasferimento dell'intero componente edilizio relativa agli $N$  strati:
\begin{equation}
    \mathbf{Z}=\mathbf{Z}_{N}\cdot\mathbf{Z}_{N-1}\cdot\ldots\cdot\mathbf{Z}_{2}\cdot\mathbf{Z}_{1} \qquad .
\end{equation}
Calcolando le matrici di trasferimento degli strati liminari come
\begin{equation}
    \mathbf{Z}_{\text{int.}} = 
    \begin{bmatrix}
        1 & -R_{\text{s,int.}} \\
        0 & 1
    \end{bmatrix}
    \qquad
    \mathbf{Z}_{\text{est.}} = 
    \begin{bmatrix}
        1 & -R_{\text{s,est.}} \\
        0 & 1
    \end{bmatrix}
\end{equation}
si può infine ottenere la matrice di trasferimento da ambiente ad ambiente:
\begin{equation}
    \mathbf{Z}_{\text{ee}} = \mathbf{Z}_{\text{est.}} \cdot \mathbf{Z} \cdot \mathbf{Z}_{\text{int.}} \qquad .
\end{equation}

Dagli elementi di questa matrice è possibile ricavare le grandezze di nostro interesse per l'analisi dinamica.
Dall'elemento {\footnotesize $12$} della matrice $\mathbf{Z_{\text{ee}}}$ si ricava la trasmittanza termica periodica $Y_{12}$ tramite il modulo della parte reale e immaginaria
\begin{equation}
    Y_{12} = \left\lvert -\frac{1}{Z_{\text{ee},12}}\right\rvert 
\end{equation} 
e il fattore di attenuazione $f$
\begin{equation}
    f = -\frac{Y_{12}}{U} \qquad .
\end{equation}
Lo sfasamento riferito al tempo iniziale si ottiene infine come:
\begin{equation}
    \Delta T = \frac{T}{2\,\pi}\arg\left(Z_{\text{ee},12}\right) + \frac{T}{2} \qquad .
\end{equation}
Dagli elementi {\footnotesize $11$} e {\footnotesize $22$} si ricavano riispettivamente l'ammettenza termica interna ed esterna 
\begin{equation}
    Y_{\text{ii}}=Y_{11} = \left\lvert -\frac{1}{Z_{\text{ee},11}}\right\rvert \qquad \qquad
    Y_{\text{ee}}=Y_{22} = \left\lvert -\frac{1}{Z_{\text{ee},22}}\right\rvert 
\end{equation} 

\begin{landscape}
        \begin{table}[H]
\centering
\resizebox{\linewidth}{!}{%
\begin{tabular}{@{}
				l
				S[table-format=1.3]
				S[table-format=4.1]
				S[table-format=4.1]
				S[table-format=2.2]
				S[table-format=1.3]
				S[table-format=1.3]
				S[table-format=2.1]
				S[table-format=1.1]
				S[table-format=4.1]
				@{}
				}
\toprule
\multicolumn{1}{c}{\multirow{3}{*}{Strati}} & \multicolumn{1}{c}{\multirow{2}{*}{Spessori}}    & \multicolumn{1}{c}{\multirow{2}{*}{Densità}}                                            & \multicolumn{1}{c}{Calore}                                            & \multicolumn{1}{c}{Massa}                                               & \multicolumn{1}{c}{Profondità di}                      & \multicolumn{1}{c}{Rapporto}   & \multicolumn{1}{c}{Capacità}                                                  & \multicolumn{1}{c}{Diffusività}                                      & \multicolumn{1}{c}{Effusività}                                                             \\
\multicolumn{1}{c}{}                        & \multicolumn{1}{c}{}                              & \multicolumn{1}{c}{}                                                  & \multicolumn{1}{c}{specifico}                                         & \multicolumn{1}{c}{superficiale}                                        & \multicolumn{1}{c}{penetrazione $\delta$} & \multicolumn{1}{c}{$\xi$}      & \multicolumn{1}{c}{termica areica}                                            & \multicolumn{1}{c}{termica}                                          & \multicolumn{1}{c}{Termica}                                                                \\
\multicolumn{1}{c}{}                        & \multicolumn{1}{c}{$\left[\SI{}{\metre}\right]$} & \multicolumn{1}{c}{$\left[\SI{}{\kilo\gram\per\metre\cubed}\right]$} & \multicolumn{1}{c}{$\left[\SI{}{\joule\per\kilo\gram\per\kelvin}\right]$} & \multicolumn{1}{c}{$\left[\SI{}{\kilo\gram\per\metre\squared}\right]$} & \multicolumn{1}{c}{$\left[\SI{}{\metre}\right]$}      & \multicolumn{1}{c}{$[-]$} & \multicolumn{1}{c}{$\left[\SI{}{\kilo\joule\per\metre\squared\per\kelvin}\right]$} & \multicolumn{1}{c}{$\left[\SI{.e-7}{\metre\squared\per\second}\right]$} & \multicolumn{1}{c}{$\left[\SI{}{\watt\second\tothe{0.5}\per\metre\squared\per\kelvin}\right]$} \\
\midrule
\end{tabular}%
}
\end{table}

\begin{flushleft}
\begin{align*}
\text{Massa superficiale totale} \, M_s &= \SI{257.9}{\kilo\gram\per\metre\squared}\\
\text{Sfasamento} \, \Delta\tau &= \SI{9.09}{\hour}\\
\text{Fattore di attenuazione} \, fd &= \SI{0.352}{}\\
\text{Trasmittanza termica periodica} \, Y_{12} &= \SI{0.088}{\watt\per\metre\squared\per\kelvin}\\
\text{Ammettanza termica interna} \, Y_{11} &= \SI{1.744}{\watt\per\metre\squared\per\kelvin}\\
\text{Ammettanza termica esterna} \, Y_{22} &= \SI{5.991}{\watt\per\metre\squared\per\kelvin}\\
\text{Capacità termica periodica interna} \, k_1 &= \SI{25.09}{\kilo\joule\per\metre\squared\per\kelvin}\\
\end{align*}
\end{flushleft}

        \clearpage
        \begin{table}[H]
\centering
\resizebox{\linewidth}{!}{%
\begin{tabular}{@{}
				l
				S[table-format=1.3]
				S[table-format=4.1]
				S[table-format=4.1]
				S[table-format=2.2]
				S[table-format=1.3]
				S[table-format=1.3]
				S[table-format=2.1]
				S[table-format=1.1]
				S[table-format=4.1]
				@{}
				}
\toprule
\multicolumn{1}{c}{\multirow{3}{*}{Strati}} & \multicolumn{1}{c}{\multirow{2}{*}{Spessori}}    & \multicolumn{1}{c}{\multirow{2}{*}{Densità}}                                            & \multicolumn{1}{c}{Calore}                                            & \multicolumn{1}{c}{Massa}                                               & \multicolumn{1}{c}{Profondità di}                      & \multicolumn{1}{c}{Rapporto}   & \multicolumn{1}{c}{Capacità}                                                  & \multicolumn{1}{c}{Diffusività}                                      & \multicolumn{1}{c}{Effusività}                                                             \\
\multicolumn{1}{c}{}                        & \multicolumn{1}{c}{}                              & \multicolumn{1}{c}{}                                                  & \multicolumn{1}{c}{specifico}                                         & \multicolumn{1}{c}{superficiale}                                        & \multicolumn{1}{c}{penetrazione $\delta$} & \multicolumn{1}{c}{$\xi$}      & \multicolumn{1}{c}{termica areica}                                            & \multicolumn{1}{c}{termica}                                          & \multicolumn{1}{c}{Termica}                                                                \\
\multicolumn{1}{c}{}                        & \multicolumn{1}{c}{$\left[\SI{}{\metre}\right]$} & \multicolumn{1}{c}{$\left[\SI{}{\kilo\gram\per\metre\cubed}\right]$} & \multicolumn{1}{c}{$\left[\SI{}{\joule\per\kilo\gram\per\kelvin}\right]$} & \multicolumn{1}{c}{$\left[\SI{}{\kilo\gram\per\metre\squared}\right]$} & \multicolumn{1}{c}{$\left[\SI{}{\metre}\right]$}      & \multicolumn{1}{c}{$[-]$} & \multicolumn{1}{c}{$\left[\SI{}{\kilo\joule\per\metre\squared\per\kelvin}\right]$} & \multicolumn{1}{c}{$\left[\SI{.e-7}{\metre\squared\per\second}\right]$} & \multicolumn{1}{c}{$\left[\SI{}{\watt\second\tothe{0.5}\per\metre\squared\per\kelvin}\right]$} \\
\midrule
              Intonaco &    0,015 &  1500,0 &           1000,0 &               22,5 &                      0,105 &  0,143 &             22,5 &                       4,00 &              948,7 \\
   Laterizio semipieno &    0,200 &  1000,0 &            840,0 &              200,0 &                      0,132 &  1,518 &            168,0 &                       6,31 &              667,2 \\
 Isolante VentirockDuo &    0,120 &    70,0 &           1030,0 &                8,4 &                      0,116 &  1,039 &              8,7 &                       4,85 &               50,2 \\
              Intonaco &    0,015 &  1800,0 &           1000,0 &               27,0 &                      0,117 &  0,128 &             27,0 &                       5,00 &             1272,8 \\
\bottomrule
\end{tabular}%
}
\end{table}

\begin{flushleft}
\begin{align*}
\text{Massa superficiale totale} \, M_s &= \SI{257.9}{\kilo\gram\per\metre\squared}\\
\text{Sfasamento} \, \Delta\tau &= \SI{9.44}{\hour}\\
\text{Fattore di attenuazione} \, fd &= \SI{0.252}{}\\
\text{Trasmittanza termica periodica} \, Y_{12} &= \SI{0.063}{\watt\per\metre\squared\per\kelvin}\\
\text{Ammettanza termica interna} \, Y_{11} &= \SI{3.925}{\watt\per\metre\squared\per\kelvin}\\
\text{Ammettanza termica esterna} \, Y_{22} &= \SI{2.157}{\watt\per\metre\squared\per\kelvin}\\
\text{Capacità termica periodica interna} \, k_1 &= \SI{54.82}{\kilo\joule\per\metre\squared\per\kelvin}\\
\end{align*}
\end{flushleft}

        \clearpage
        \begin{table}[H]
\centering
\resizebox{\linewidth}{!}{%
\begin{tabular}{@{}
				l
				S[table-format=1.3]
				S[table-format=4.1]
				S[table-format=4.1]
				S[table-format=2.2]
				S[table-format=1.3]
				S[table-format=1.3]
				S[table-format=2.1]
				S[table-format=1.1]
				S[table-format=4.1]
				@{}
				}
\toprule
\multicolumn{1}{c}{\multirow{3}{*}{Strati}} & \multicolumn{1}{c}{\multirow{2}{*}{Spessori}}    & \multicolumn{1}{c}{\multirow{2}{*}{Densità}}                                            & \multicolumn{1}{c}{Calore}                                            & \multicolumn{1}{c}{Massa}                                               & \multicolumn{1}{c}{Profondità di}                      & \multicolumn{1}{c}{Rapporto}   & \multicolumn{1}{c}{Capacità}                                                  & \multicolumn{1}{c}{Diffusività}                                      & \multicolumn{1}{c}{Effusività}                                                             \\
\multicolumn{1}{c}{}                        & \multicolumn{1}{c}{}                              & \multicolumn{1}{c}{}                                                  & \multicolumn{1}{c}{specifico}                                         & \multicolumn{1}{c}{superficiale}                                        & \multicolumn{1}{c}{penetrazione $\delta$} & \multicolumn{1}{c}{$\xi$}      & \multicolumn{1}{c}{termica areica}                                            & \multicolumn{1}{c}{termica}                                          & \multicolumn{1}{c}{Termica}                                                                \\
\multicolumn{1}{c}{}                        & \multicolumn{1}{c}{$\left[\SI{}{\metre}\right]$} & \multicolumn{1}{c}{$\left[\SI{}{\kilo\gram\per\metre\cubed}\right]$} & \multicolumn{1}{c}{$\left[\SI{}{\joule\per\kilo\gram\per\kelvin}\right]$} & \multicolumn{1}{c}{$\left[\SI{}{\kilo\gram\per\metre\squared}\right]$} & \multicolumn{1}{c}{$\left[\SI{}{\metre}\right]$}      & \multicolumn{1}{c}{$[-]$} & \multicolumn{1}{c}{$\left[\SI{}{\kilo\joule\per\metre\squared\per\kelvin}\right]$} & \multicolumn{1}{c}{$\left[\SI{.e-7}{\metre\squared\per\second}\right]$} & \multicolumn{1}{c}{$\left[\SI{}{\watt\second\tothe{0.5}\per\metre\squared\per\kelvin}\right]$} \\
\midrule
             Gessofibra &    0,015 &  1150,0 &           1100,0 &              17,25 &                      0,068 &  0,222 &             19,0 &                       1,66 &              515,4 \\
              X-LAM KLH &    0,096 &   500,0 &           1600,0 &              48,00 &                      0,067 &  1,436 &             76,8 &                       1,62 &              322,5 \\
 Isolante bassa densità &    0,120 &    50,0 &           2100,0 &               6,00 &                      0,100 &  1,203 &             12,6 &                       3,62 &               63,2 \\
         Intonaco calce &    0,015 &  1800,0 &           1000,0 &              27,00 &                      0,117 &  0,128 &             27,0 &                       5,00 &             1272,8 \\
\bottomrule
\end{tabular}%
}
\end{table}

\begin{flushleft}
\begin{align*}
\text{Massa superficiale totale} \, M_s &= \SI{98.25}{\kilo\gram\per\metre\squared}\\
\text{Sfasamento} \, \Delta\tau &= \SI{9.33}{\hour}\\
\text{Fattore di attenuazione} \, fd &= \SI{0.334}{}\\
\text{Trasmittanza termica periodica} \, Y_{12} &= \SI{0.08}{\watt\per\metre\squared\per\kelvin}\\
\text{Ammettanza termica interna} \, Y_{11} &= \SI{2.651}{\watt\per\metre\squared\per\kelvin}\\
\text{Ammettanza termica esterna} \, Y_{22} &= \SI{2.257}{\watt\per\metre\squared\per\kelvin}\\
\text{Capacità termica periodica interna} \, k_1 &= \SI{37.56}{\kilo\joule\per\metre\squared\per\kelvin}\\
\end{align*}
\end{flushleft}

        \clearpage
        \begin{table}
\centering
\caption{Perete in X-LAM con isolante bassa densità lana di roccia}
\begin{tabular}{lrrrrrr}
\toprule
                Strati & Spessori & Densità & Calore specifico & Massa superficiale & Profondità di Penetrazione &     xi \\
\midrule
            Gessofibra &    0,013 &  1150,0 &           1100,0 &              14,95 &                      0,068 &  0,192 \\
             X-LAM KLH &    0,096 &   500,0 &           1600,0 &              48,00 &                      0,067 &  1,436 \\
 Isolante ventirockduo &    0,105 &    70,0 &           1030,0 &               7,35 &                      0,116 &  0,909 \\
        Intonaco calce &    0,015 &  1800,0 &           1000,0 &              27,00 &                      0,117 &  0,128 \\
\bottomrule
\end{tabular}
\end{table}

\begin{flushleft}
\begin{align*}
\text{Massa superficiale totale} &= \SI{97.3}{\kilo\gram}\\
\text{Sfasamento} &= \SI{8.34}{\hour}\\
\text{Attenuazione} &= \SI{239.221}{}
\end{align*}
\end{flushleft}

        \clearpage
        \begin{table}[H]
\centering
\resizebox{\linewidth}{!}{%
\begin{tabular}{@{}
				l
				S[table-format=1.3]
				S[table-format=4.1]
				S[table-format=4.1]
				S[table-format=2.2]
				S[table-format=1.3]
				S[table-format=1.3]
				S[table-format=2.1]
				S[table-format=1.1]
				S[table-format=4.1]
				@{}
				}
\toprule
\multicolumn{1}{c}{\multirow{3}{*}{Strati}} & \multicolumn{1}{c}{\multirow{2}{*}{Spessori}}    & \multicolumn{1}{c}{\multirow{2}{*}{Densità}}                                            & \multicolumn{1}{c}{Calore}                                            & \multicolumn{1}{c}{Massa}                                               & \multicolumn{1}{c}{Profondità di}                      & \multicolumn{1}{c}{Rapporto}   & \multicolumn{1}{c}{Capacità}                                                  & \multicolumn{1}{c}{Diffusività}                                      & \multicolumn{1}{c}{Effusività}                                                             \\
\multicolumn{1}{c}{}                        & \multicolumn{1}{c}{}                              & \multicolumn{1}{c}{}                                                  & \multicolumn{1}{c}{specifico}                                         & \multicolumn{1}{c}{superficiale}                                        & \multicolumn{1}{c}{penetrazione $\delta$} & \multicolumn{1}{c}{$\xi$}      & \multicolumn{1}{c}{termica areica}                                            & \multicolumn{1}{c}{termica}                                          & \multicolumn{1}{c}{Termica}                                                                \\
\multicolumn{1}{c}{}                        & \multicolumn{1}{c}{$\left[\SI{}{\metre}\right]$} & \multicolumn{1}{c}{$\left[\SI{}{\kilo\gram\per\metre\cubed}\right]$} & \multicolumn{1}{c}{$\left[\SI{}{\joule\per\kilo\gram\per\kelvin}\right]$} & \multicolumn{1}{c}{$\left[\SI{}{\kilo\gram\per\metre\squared}\right]$} & \multicolumn{1}{c}{$\left[\SI{}{\metre}\right]$}      & \multicolumn{1}{c}{$[-]$} & \multicolumn{1}{c}{$\left[\SI{}{\kilo\joule\per\metre\squared\per\kelvin}\right]$} & \multicolumn{1}{c}{$\left[\SI{.e-7}{\metre\squared\per\second}\right]$} & \multicolumn{1}{c}{$\left[\SI{}{\watt\second\tothe{0.5}\per\metre\squared\per\kelvin}\right]$} \\
\midrule
                                Gessofibra &    0,015 &  1150,0 &           1100,0 &              17,25 &                      0,068 &  0,222 &             19,0 &                       1,66 &              515,4 \\
                                 X-LAM KLH &    0,096 &   500,0 &           1600,0 &              48,00 &                      0,067 &  1,436 &             76,8 &                       1,62 &              322,5 \\
 Isolante alta densità Naturalia Diffuterm &    0,130 &   190,0 &           2100,0 &              24,70 &                      0,054 &  2,388 &             51,9 &                       1,08 &              131,0 \\
                            Intonaco calce &    0,015 &  1800,0 &           1000,0 &              27,00 &                      0,117 &  0,128 &             27,0 &                       5,00 &             1272,8 \\
\bottomrule
\end{tabular}%
}
\end{table}

\begin{flushleft}
\begin{align*}
\text{Massa superficiale totale} \, M_s &= \SI{116.95}{\kilo\gram\per\metre\squared}\\
\text{Sfasamento} \, \Delta\tau &= \SI{13.78}{\hour}\\
\text{Fattore di attenuazione} \, fd &= \SI{0.177}{}\\
\text{Trasmittanza termica periodica} \, Y_{12} &= \SI{0.044}{\watt\per\metre\squared\per\kelvin}\\
\text{Ammettanza termica interna} \, Y_{11} &= \SI{2.613}{\watt\per\metre\squared\per\kelvin}\\
\text{Ammettanza termica esterna} \, Y_{22} &= \SI{2.73}{\watt\per\metre\squared\per\kelvin}\\
\text{Capacità termica periodica interna} \, k_1 &= \SI{36.21}{\kilo\joule\per\metre\squared\per\kelvin}\\
\end{align*}
\end{flushleft}

        \clearpage
        % This file was created by tikzplotlib v0.9.4.
\begin{tikzpicture}

\definecolor{color0}{rgb}{0.12156862745098,0.466666666666667,0.705882352941177}
\definecolor{color1}{rgb}{1,0.498039215686275,0.0549019607843137}
\definecolor{color2}{rgb}{0.172549019607843,0.627450980392157,0.172549019607843}
\definecolor{color3}{rgb}{0.83921568627451,0.152941176470588,0.156862745098039}
\definecolor{color4}{rgb}{0.580392156862745,0.403921568627451,0.741176470588235}
\definecolor{color5}{rgb}{0.549019607843137,0.337254901960784,0.294117647058824}
\definecolor{color6}{rgb}{0.890196078431372,0.466666666666667,0.76078431372549}

\begin{axis}[
height=15cm,
legend cell align={left},
legend style={fill opacity=0.8, draw opacity=1, text opacity=1, draw=white!80!black},
tick align=outside,
tick pos=left,
width=\linewidth,
x grid style={white!69.0196078431373!black},
xmin=-0.783714285714286, xmax=7.78371428571429,
xtick style={color=black},
xtick={-1,0,1,2,3,4,5,6,7,8},
xticklabels={,0a,0b,1a,1b,2a,3a,3b,3c,},
y grid style={white!69.0196078431373!black},
ymin=0, ymax=28.3007380178985,
ytick style={color=black}
]
\draw[draw=none,fill=color0] (axis cs:-0.394285714285714,0) rectangle (axis cs:-0.291428571428571,1);
\draw[draw=none,fill=color0] (axis cs:0.605714285714286,0) rectangle (axis cs:0.708571428571429,0.12);
\draw[draw=none,fill=color0] (axis cs:1.60571428571429,0) rectangle (axis cs:1.70857142857143,0.35);
\draw[draw=none,fill=color0] (axis cs:2.60571428571429,0) rectangle (axis cs:2.70857142857143,0.35);
\draw[draw=none,fill=color0] (axis cs:3.60571428571429,0) rectangle (axis cs:3.70857142857143,0.34);
\draw[draw=none,fill=color0] (axis cs:4.60571428571429,0) rectangle (axis cs:4.70857142857143,0.246);
\draw[draw=none,fill=color0] (axis cs:5.60571428571429,0) rectangle (axis cs:5.70857142857143,0.236);
\draw[draw=none,fill=color0] (axis cs:6.60571428571429,0) rectangle (axis cs:6.70857142857143,0.256);
\draw[draw=none,fill=color1] (axis cs:-0.28,0) rectangle (axis cs:-0.177142857142857,1.45205128205128);
\draw[draw=none,fill=color1] (axis cs:0.72,0) rectangle (axis cs:0.822857142857143,3.59857142857143);
\draw[draw=none,fill=color1] (axis cs:1.72,0) rectangle (axis cs:1.82285714285714,4.01759658580413);
\draw[draw=none,fill=color1] (axis cs:2.72,0) rectangle (axis cs:2.82285714285714,4.01759658580413);
\draw[draw=none,fill=color1] (axis cs:3.72,0) rectangle (axis cs:3.82285714285714,4.22408902691511);
\draw[draw=none,fill=color1] (axis cs:4.72,0) rectangle (axis cs:4.82285714285714,4.15445151339888);
\draw[draw=none,fill=color1] (axis cs:5.72,0) rectangle (axis cs:5.82285714285714,4.13941391941392);
\draw[draw=none,fill=color1] (axis cs:6.72,0) rectangle (axis cs:6.82285714285714,4.01981259051027);
\draw[draw=none,fill=color2] (axis cs:-0.165714285714286,0) rectangle (axis cs:-0.0628571428571428,17);
\draw[draw=none,fill=color2] (axis cs:0.834285714285714,0) rectangle (axis cs:0.937142857142857,0.084);
\draw[draw=none,fill=color2] (axis cs:1.83428571428571,0) rectangle (axis cs:1.93714285714286,2.579);
\draw[draw=none,fill=color2] (axis cs:2.83428571428571,0) rectangle (axis cs:2.93714285714286,2.579);
\draw[draw=none,fill=color2] (axis cs:3.83428571428571,0) rectangle (axis cs:3.93714285714286,2.292);
\draw[draw=none,fill=color2] (axis cs:4.83428571428571,0) rectangle (axis cs:4.93714285714286,0.9825);
\draw[draw=none,fill=color2] (axis cs:5.83428571428571,0) rectangle (axis cs:5.93714285714286,0.9995);
\draw[draw=none,fill=color2] (axis cs:6.83428571428571,0) rectangle (axis cs:6.93714285714286,1.1695);
\draw[draw=none,fill=color3] (axis cs:-0.0514285714285714,0) rectangle (axis cs:0.0514285714285714,0.12657997176382);
\draw[draw=none,fill=color3] (axis cs:0.948571428571429,0) rectangle (axis cs:1.05142857142857,26.95308382657);
\draw[draw=none,fill=color3] (axis cs:1.94857142857143,0) rectangle (axis cs:2.05142857142857,8.7574995923251);
\draw[draw=none,fill=color3] (axis cs:2.94857142857143,0) rectangle (axis cs:3.05142857142857,6.27975484198694);
\draw[draw=none,fill=color3] (axis cs:3.94857142857143,0) rectangle (axis cs:4.05142857142857,5.29986171570837);
\draw[draw=none,fill=color3] (axis cs:4.94857142857143,0) rectangle (axis cs:5.05142857142857,8.04843003171711);
\draw[draw=none,fill=color3] (axis cs:5.94857142857143,0) rectangle (axis cs:6.05142857142857,8.5884800548406);
\draw[draw=none,fill=color3] (axis cs:6.94857142857143,0) rectangle (axis cs:7.05142857142857,4.39742417434425);
\draw[draw=none,fill=color4] (axis cs:0.0628571428571429,0) rectangle (axis cs:0.165714285714286,8.43432943732685);
\draw[draw=none,fill=color4] (axis cs:1.06285714285714,0) rectangle (axis cs:1.16571428571429,1.48909024306196);
\draw[draw=none,fill=color4] (axis cs:2.06285714285714,0) rectangle (axis cs:2.16571428571429,9.08756910878924);
\draw[draw=none,fill=color4] (axis cs:3.06285714285714,0) rectangle (axis cs:3.16571428571429,9.4429593990121);
\draw[draw=none,fill=color4] (axis cs:4.06285714285714,0) rectangle (axis cs:4.16571428571429,10.9678887890606);
\draw[draw=none,fill=color4] (axis cs:5.06285714285714,0) rectangle (axis cs:5.16571428571429,9.33067092051128);
\draw[draw=none,fill=color4] (axis cs:6.06285714285714,0) rectangle (axis cs:6.16571428571429,8.5813220859026);
\draw[draw=none,fill=color4] (axis cs:7.06285714285714,0) rectangle (axis cs:7.16571428571429,13.7756666736772);
\draw[draw=none,fill=color5] (axis cs:0.177142857142857,0) rectangle (axis cs:0.28,0.0183800610281671);
\draw[draw=none,fill=color5] (axis cs:1.17714285714286,0) rectangle (axis cs:1.28,9.69925973701853);
\draw[draw=none,fill=color5] (axis cs:2.17714285714286,0) rectangle (axis cs:2.28,3.51841004623064);
\draw[draw=none,fill=color5] (axis cs:3.17714285714286,0) rectangle (axis cs:3.28,2.52295216128537);
\draw[draw=none,fill=color5] (axis cs:4.17714285714286,0) rectangle (axis cs:4.28,2.23870877174912);
\draw[draw=none,fill=color5] (axis cs:5.17714285714286,0) rectangle (axis cs:5.28,3.34368123257522);
\draw[draw=none,fill=color5] (axis cs:6.17714285714286,0) rectangle (axis cs:6.28,3.5551273885616);
\draw[draw=none,fill=color5] (axis cs:7.17714285714286,0) rectangle (axis cs:7.28,1.76768210618432);
\draw[draw=none,fill=color6] (axis cs:0.291428571428571,0) rectangle (axis cs:0.394285714285714,6.32800639964363);
\draw[draw=none,fill=color6] (axis cs:1.29142857142857,0) rectangle (axis cs:1.39428571428571,0.410482190163323);
\draw[draw=none,fill=color6] (axis cs:2.29142857142857,0) rectangle (axis cs:2.39428571428571,2.5094786899098);
\draw[draw=none,fill=color6] (axis cs:3.29142857142857,0) rectangle (axis cs:3.39428571428571,5.48209087513);
\draw[draw=none,fill=color6] (axis cs:4.29142857142857,0) rectangle (axis cs:4.39428571428571,2.43814076713028);
\draw[draw=none,fill=color6] (axis cs:5.29142857142857,0) rectangle (axis cs:5.39428571428571,3.75582315594814);
\draw[draw=none,fill=color6] (axis cs:6.29142857142857,0) rectangle (axis cs:6.39428571428571,3.77092834381981);
\draw[draw=none,fill=color6] (axis cs:7.29142857142857,0) rectangle (axis cs:7.39428571428571,3.62094494519863);
\draw (axis cs:-0.342857142857143,1) ++(0pt,3pt) node[
  scale=0.5,
  anchor=south,
  text=black,
  rotate=0.0
]{1.0};
\draw (axis cs:0.657142857142857,0.12) ++(0pt,3pt) node[
  scale=0.5,
  anchor=south,
  text=black,
  rotate=0.0
]{0.12};
\draw (axis cs:1.65714285714286,0.35) ++(0pt,3pt) node[
  scale=0.5,
  anchor=south,
  text=black,
  rotate=0.0
]{0.35};
\draw (axis cs:2.65714285714286,0.35) ++(0pt,3pt) node[
  scale=0.5,
  anchor=south,
  text=black,
  rotate=0.0
]{0.35};
\draw (axis cs:3.65714285714286,0.34) ++(0pt,3pt) node[
  scale=0.5,
  anchor=south,
  text=black,
  rotate=0.0
]{0.34};
\draw (axis cs:4.65714285714286,0.25) ++(0pt,3pt) node[
  scale=0.5,
  anchor=south,
  text=black,
  rotate=0.0
]{0.25};
\draw (axis cs:5.65714285714286,0.24) ++(0pt,3pt) node[
  scale=0.5,
  anchor=south,
  text=black,
  rotate=0.0
]{0.24};
\draw (axis cs:6.65714285714286,0.26) ++(0pt,3pt) node[
  scale=0.5,
  anchor=south,
  text=black,
  rotate=0.0
]{0.26};
\draw (axis cs:-0.228571428571429,1.45) ++(0pt,3pt) node[
  scale=0.5,
  anchor=south,
  text=black,
  rotate=0.0
]{1.45};
\draw (axis cs:0.771428571428571,3.6) ++(0pt,3pt) node[
  scale=0.5,
  anchor=south,
  text=black,
  rotate=0.0
]{3.6};
\draw (axis cs:1.77142857142857,4.02) ++(0pt,3pt) node[
  scale=0.5,
  anchor=south,
  text=black,
  rotate=0.0
]{4.02};
\draw (axis cs:2.77142857142857,4.02) ++(0pt,3pt) node[
  scale=0.5,
  anchor=south,
  text=black,
  rotate=0.0
]{4.02};
\draw (axis cs:3.77142857142857,4.22) ++(0pt,3pt) node[
  scale=0.5,
  anchor=south,
  text=black,
  rotate=0.0
]{4.22};
\draw (axis cs:4.77142857142857,4.15) ++(0pt,3pt) node[
  scale=0.5,
  anchor=south,
  text=black,
  rotate=0.0
]{4.15};
\draw (axis cs:5.77142857142857,4.14) ++(0pt,3pt) node[
  scale=0.5,
  anchor=south,
  text=black,
  rotate=0.0
]{4.14};
\draw (axis cs:6.77142857142857,4.02) ++(0pt,3pt) node[
  scale=0.5,
  anchor=south,
  text=black,
  rotate=0.0
]{4.02};
\draw (axis cs:-0.114285714285714,17) ++(0pt,3pt) node[
  scale=0.5,
  anchor=south,
  text=black,
  rotate=0.0
]{17.0};
\draw (axis cs:0.885714285714286,0.08) ++(0pt,3pt) node[
  scale=0.5,
  anchor=south,
  text=black,
  rotate=0.0
]{0.08};
\draw (axis cs:1.88571428571429,2.58) ++(0pt,3pt) node[
  scale=0.5,
  anchor=south,
  text=black,
  rotate=0.0
]{2.58};
\draw (axis cs:2.88571428571429,2.58) ++(0pt,3pt) node[
  scale=0.5,
  anchor=south,
  text=black,
  rotate=0.0
]{2.58};
\draw (axis cs:3.88571428571429,2.29) ++(0pt,3pt) node[
  scale=0.5,
  anchor=south,
  text=black,
  rotate=0.0
]{2.29};
\draw (axis cs:4.88571428571429,0.98) ++(0pt,3pt) node[
  scale=0.5,
  anchor=south,
  text=black,
  rotate=0.0
]{0.98};
\draw (axis cs:5.88571428571429,1) ++(0pt,3pt) node[
  scale=0.5,
  anchor=south,
  text=black,
  rotate=0.0
]{1.0};
\draw (axis cs:6.88571428571429,1.17) ++(0pt,3pt) node[
  scale=0.5,
  anchor=south,
  text=black,
  rotate=0.0
]{1.17};
\draw (axis cs:0,0.13) ++(0pt,3pt) node[
  scale=0.5,
  anchor=south,
  text=black,
  rotate=0.0
]{0.13};
\draw (axis cs:1,26.95) ++(0pt,3pt) node[
  scale=0.5,
  anchor=south,
  text=black,
  rotate=0.0
]{26.95};
\draw (axis cs:2,8.76) ++(0pt,3pt) node[
  scale=0.5,
  anchor=south,
  text=black,
  rotate=0.0
]{8.76};
\draw (axis cs:3,6.28) ++(0pt,3pt) node[
  scale=0.5,
  anchor=south,
  text=black,
  rotate=0.0
]{6.28};
\draw (axis cs:4,5.3) ++(0pt,3pt) node[
  scale=0.5,
  anchor=south,
  text=black,
  rotate=0.0
]{5.3};
\draw (axis cs:5,8.05) ++(0pt,3pt) node[
  scale=0.5,
  anchor=south,
  text=black,
  rotate=0.0
]{8.05};
\draw (axis cs:6,8.59) ++(0pt,3pt) node[
  scale=0.5,
  anchor=south,
  text=black,
  rotate=0.0
]{8.59};
\draw (axis cs:7,4.4) ++(0pt,3pt) node[
  scale=0.5,
  anchor=south,
  text=black,
  rotate=0.0
]{4.4};
\draw (axis cs:0.114285714285714,8.43) ++(0pt,3pt) node[
  scale=0.5,
  anchor=south,
  text=black,
  rotate=0.0
]{8.43};
\draw (axis cs:1.11428571428571,1.49) ++(0pt,3pt) node[
  scale=0.5,
  anchor=south,
  text=black,
  rotate=0.0
]{1.49};
\draw (axis cs:2.11428571428571,9.09) ++(0pt,3pt) node[
  scale=0.5,
  anchor=south,
  text=black,
  rotate=0.0
]{9.09};
\draw (axis cs:3.11428571428571,9.44) ++(0pt,3pt) node[
  scale=0.5,
  anchor=south,
  text=black,
  rotate=0.0
]{9.44};
\draw (axis cs:4.11428571428571,10.97) ++(0pt,3pt) node[
  scale=0.5,
  anchor=south,
  text=black,
  rotate=0.0
]{10.97};
\draw (axis cs:5.11428571428571,9.33) ++(0pt,3pt) node[
  scale=0.5,
  anchor=south,
  text=black,
  rotate=0.0
]{9.33};
\draw (axis cs:6.11428571428571,8.58) ++(0pt,3pt) node[
  scale=0.5,
  anchor=south,
  text=black,
  rotate=0.0
]{8.58};
\draw (axis cs:7.11428571428571,13.78) ++(0pt,3pt) node[
  scale=0.5,
  anchor=south,
  text=black,
  rotate=0.0
]{13.78};
\draw (axis cs:0.228571428571429,0.02) ++(0pt,3pt) node[
  scale=0.5,
  anchor=south,
  text=black,
  rotate=0.0
]{0.02};
\draw (axis cs:1.22857142857143,9.7) ++(0pt,3pt) node[
  scale=0.5,
  anchor=south,
  text=black,
  rotate=0.0
]{9.7};
\draw (axis cs:2.22857142857143,3.52) ++(0pt,3pt) node[
  scale=0.5,
  anchor=south,
  text=black,
  rotate=0.0
]{3.52};
\draw (axis cs:3.22857142857143,2.52) ++(0pt,3pt) node[
  scale=0.5,
  anchor=south,
  text=black,
  rotate=0.0
]{2.52};
\draw (axis cs:4.22857142857143,2.24) ++(0pt,3pt) node[
  scale=0.5,
  anchor=south,
  text=black,
  rotate=0.0
]{2.24};
\draw (axis cs:5.22857142857143,3.34) ++(0pt,3pt) node[
  scale=0.5,
  anchor=south,
  text=black,
  rotate=0.0
]{3.34};
\draw (axis cs:6.22857142857143,3.56) ++(0pt,3pt) node[
  scale=0.5,
  anchor=south,
  text=black,
  rotate=0.0
]{3.56};
\draw (axis cs:7.22857142857143,1.77) ++(0pt,3pt) node[
  scale=0.5,
  anchor=south,
  text=black,
  rotate=0.0
]{1.77};
\draw (axis cs:0.342857142857143,6.33) ++(0pt,3pt) node[
  scale=0.5,
  anchor=south,
  text=black,
  rotate=0.0
]{6.33};
\draw (axis cs:1.34285714285714,0.41) ++(0pt,3pt) node[
  scale=0.5,
  anchor=south,
  text=black,
  rotate=0.0
]{0.41};
\draw (axis cs:2.34285714285714,2.51) ++(0pt,3pt) node[
  scale=0.5,
  anchor=south,
  text=black,
  rotate=0.0
]{2.51};
\draw (axis cs:3.34285714285714,5.48) ++(0pt,3pt) node[
  scale=0.5,
  anchor=south,
  text=black,
  rotate=0.0
]{5.48};
\draw (axis cs:4.34285714285714,2.44) ++(0pt,3pt) node[
  scale=0.5,
  anchor=south,
  text=black,
  rotate=0.0
]{2.44};
\draw (axis cs:5.34285714285714,3.76) ++(0pt,3pt) node[
  scale=0.5,
  anchor=south,
  text=black,
  rotate=0.0
]{3.76};
\draw (axis cs:6.34285714285714,3.77) ++(0pt,3pt) node[
  scale=0.5,
  anchor=south,
  text=black,
  rotate=0.0
]{3.77};
\draw (axis cs:7.34285714285714,3.62) ++(0pt,3pt) node[
  scale=0.5,
  anchor=south,
  text=black,
  rotate=0.0
]{3.62};
\end{axis}

\end{tikzpicture}

\end{landscape}


\end{document}

